\section{Tomberg’s Labor}

Valentin Tomberg (as any reader of this blog by now knows) plays a not insignificant role in linking Christianity with esotericism, even paganism (of the older kind). It would be easy to just go through and cherry pick random quotes from Tomberg: this would actually turn out much better than one might imagine, since Tomberg manages to stuff every sentence with profound meaning and an inner direction which helps even isolated sentences maintain their context.

There is a particular section of text which helps to elucidate his purpose more sharply and pointedly than is his wont, as Tomberg usually manages to hide in plain sight. On my first reading of the text, I apparently gleaned much less than I missed, and am ashamed to say that this only jumped out on a second reading.

\begin{quotex}
The macrocosmic sphere of paradise (St Paul's third heaven) and the microcosmic layer of Eden are the \emph{initia} (beginning) to which one is initiated in the macrocosmic initiation as well as in the microcosmic initiation. Ecstasy to the heights beyond one's self and entasy into the depths within one's self lead to knowledge of the same fundamental truth. Christian esotericism unites these two methods of initiation.
\end{quotex}


There is his entire program in a nutshell (incidentally we see that the Christian method is more about the union of the \emph{modes}, rather than emphasis on the division between the two \emph{techniques}). He is writing a secular ``John'' for a modern age. This is very like Goethe re-casting Job in \emph{Faust}, except thankfully Tomberg actually claims to know what he is doing, and seems to back this up. In other words, Valentin Tomberg is actually claiming that a Christian can mingle the exoteric and esoteric modes of knowing in their own person, without (and this is the important point) compromising or losing either one individually.

If this claim reminds you of the claims regarding Christ and the \emph{homo ouison} with or without iotas (made by successive councils\footnote{\url{http://en.wikipedia.org/wiki/First_Council_of_Nicaea}}) it should, for that is actually his aim. The Christian is a ``little Christ'', and Tomberg wishes to perfect Nature (exoteric religion) with Super-Nature (esoteric meaning), without however losing either the Self or the ``Self beyond the Selves''. Jesus is the man-God, and Christians are also called to be fully God and truly man. He is not splitting hairs, he is describing or coming to terms with, something \emph{he had already fully experienced}.


The figure of the Emperor is the ``truest man'' (King David), and the hermit (prophet) is the most like God (fully God), remarks Tomberg in his chapter on the Pope. The pope (or priest) actually balances earth with heaven. So an invisible principle of union is here shown to be concretely uniting two seeming opposites. Here we have Soloviev's ``Ideas'', with the veil lift by for a moment, Tomberg explaining the Arcanum that shows even the invisible worlds correspond to Law rather than Chaos. For here it is not a matter of seeing contrasts, and is never a matter of contrasts. John Michael Greer brilliantly points out\footnote{\url{http://thearchdruidreport.blogspot.com/2014/03/american-delusionalism-or-why-history.html}} that focusing on differences actually obscures them (this is because, as Iamblichus teaches, you have to start with One, not Two). If you focus on similarities, the differences are seen truly and starkly for what they are.

\begin{quotex}
Any present or future set of events, however unique it may be in terms of the fine details, has points of similarity with events in the past, and those points of similarity allow the past events to be taken as a guide to the present and future. This works best if you've got a series of past events, as different from each other as any one of them is from the present or future situation you're trying to predict; if you can find common patterns in the whole range of past parallels, it's usually a safe bet that the same pattern will recur again. Any time you approach a present or future event, then, you have two choices: you can look for the features that event has in common with other events, despite the differences of detail, or you can focus on the differences and ignore the common features. The first of those choices, it's worth noting, allows you to consider both the similarities and the differences. Once you've got the common pattern, it then becomes possible to modify it as needed to take into account the special characteristics of the situation you're trying to understand or predict: to notice, for example, that the dark age that will follow our civilization will have to contend with nuclear and chemical pollution on top of the more ordinary consequences of decline and fall.If you start from the assumption that the event you're trying to predict is unlike anything that's ever happened before, though, you've thrown out your chance of perceiving the common pattern. What happens instead, with motononous regularity, is that pop-culture narratives such as the sudden overnight collapse beloved of Hollywood screenplay writers smuggle themselves into the picture, and cement themselves in place with the help of confirmation bias. The result is the endless recycling of repeatedly failed predictions that plays so central a role in the collective imagination of our time, and has helped so many people blind themselves to the unwelcome future closing in on us.
\end{quotex}


So the ``hermetic'' method is actually just \emph{sanity}, applied in depth and height, of the common-sense-rule of having a unity-principle and not forcing apples to be seen as oranges. If one sees sleep follow work, then awakening follow sleep in Nature, one can assume (by analogy) that Death follows Life, and then Life supersedes Death. This is the affirmation of the Great Chain of Being\footnote{\url{http://en.wikipedia.org/wiki/Great_chain_of_being}}, the ``Book of Nature'', the Analogia Entis.

So (too) we can say that it is the inner ``priest'' in man which reconciles the manly-earthly King with the God-mad prophet, in the dark of night, before the black turns to red, white, then gold at the dawning of the Sun. There is always a balance, and always a ``higher level''. This is because God is always God: ``higher up, and further in''. As Tomberg says, choose spiritual death, and choose hell: choose Life, and you have chosen God. It's so simple, children can do it, and so complex, that even seraphim yearn to be able to see clearly.

This occurs within Christian esotericism, and within the Christian esotericist, who privileges neither the \emph{badge} of Christianity nor the \emph{heritage} or \emph{technique} of ancient practices \emph{which are given new life through baptism}. Does anyone really suppose that it would be ``spiritually exciting'' to live under the rule of the tutelary stars in the same sense that the heathen did? So that an ill omen meant almost certain death? The heathen are a picture of our natural man, sunk in darkness. Paradoxically, it is the man who learns that he is under the rule of stars who then begins to escape that rule of the stars, and the wise man becometh free. This is why astrology is as much an art as a science, and why the Middle Ages baptized it, but \emph{did not use it to replace} the mass and the organized Christian religion. In any event, men like Pythagoras were not subject to the stars in the same sense that the local god-fearing goat-herders were. He learned to judge the angels, \& it was this that made him free. So the ``Christian'' sees that while there is ``progress'' (and man learns to be ``free'' of powers), yet these things were altogether written for our sakes\footnote{\url{http://biblehub.com/1_corinthians/9-10.htm}}, because we have the same process to over go in miniature. So, in a sense we are more free, in a sense we are less: welcome to the Kali Yuga. The point is to find God.

In the same manner, it is the true ``Christian'' who both grasps intuitively the power of the antique Tradition, \& yet sees it living through its transformations in Christianity and onto into and despite the Kali Yuga; he or she it is who can keep the contradictions together because he approaches from a standpoint of \emph{faith}, rather than doubt, and thus \emph{doesn't sink into experience} through the moving, opening, readying of wavering doubt (like Eve before the apple). This is the man who can keep \emph{entasy} (unpeeling one's soul layers) and
\emph{ecstasy} (voyaging into God) from dissolving together into the slime of muddle and confusion that tends to stamp those trying to be ``spiritual'' in the Dark Ages, using the one to heighten and enhance the other. The same wind wrecks one vessel, and lifts another one over the waves. The same wind drives down some birds, and causes the hawk to soar. Jesus the God-man is the true \emph{in hoc signo vinces}\footnote{\url{http://en.wikipedia.org/wiki/In\_hoc\_signo\_vinces}} for the esotericist or exotericist, since both are valid, and both ``mean'' the same thing. This is one of the meanings of the hypostatic union: that there is fully each, without loss of either.

Faith precedes experience because sense experience alone takes too long, and tends to lose one in a ``dark wood''. Additionally, ``faith'' is actually a divine principle which begins to work before it is either deserved or understood. It thus appears ``blind'' to the blind, and ``powerless'' to the powerless. It is actually the power and potency of God, Who is willing to act \emph{long before man is worthy of it}. Thus, it is power, because man (at that juncture) has no power. In this way, man is invited to cooperate with God, and that is the price of ascent and the meaning of the night time, which sees great growth. \emph{But this merely accords with the Nature principle of paganism}: a seed has to first fall into the ground and die, before it can have the life of the tree. Likewise, God's power is so great, it first has to die, because it has to ``make a little space'' for man by withdrawing, or else man could not be. God is so powerful, He can afford to suffer in silence. God alone can afford it. Here is the meaning both of \emph{kenosis}\footnote{\url{http://en.wikipedia.org/wiki/Kenosis}} and of grace.

We will be looking at Tomberg's journey, through the eyes of the \emph{Meditations} (and using some of the notes sent out on the mailing list, not to mention Cologero's observations and thoughts), however, it will primarily be an attempt to unfold or unpack some of his meaning, so that it stands out more clearly, so that others can ``meditate'' (which was his whole purpose anyway).


\flrightit{Posted on 2014-04-11 by Logres}

